% ============================================================
\section{Preliminaries}
\label{sec:preliminaries}
% ============================================================

This section recalls standard regular expressions, finite automata,
and the classical position automaton.  These notions provide the
baseline against which the lookahead-aware construction is defined.

\paragraph{Notation.}
We use \(a,b,\sigma\) for alphabet symbols, \(r,s\) for regular expressions, and \(u,v,w\) for words. We reserve \(p,q\) for marked positions in a regular expression and \(i,j,k\) for boundaries in an input word.  This distinction is useful because the position automaton refers to syntactic positions in a regular expression, whereas lookahead is evaluated at boundaries of the input word.

 % We write $\mathsf{T}$ and $\mathsf{F}$ for the truth and falsehood constants, respectively.


% ------------------------------------------------------------
\subsection{Regular Expressions and Finite Automata}
\label{sec:re-automata}
% ------------------------------------------------------------

Let $\Sigma$ be a non-empty finite set of symbols, called the alphabet. The set of standard regular expressions over $\Sigma$, denoted by $\mathrm{RE}$, is generated by
\begin{equation}
r,s ::= 
\emptyset
\mid \varepsilon
\mid a
\mid r+s
\mid r\cdot s
\mid r^*,
\qquad a\in\Sigma.
\label{equation: grammar for standard regex}
\end{equation}
For notational convenience, we omit the concatenation operator $\cdot$ and write $rs$ for $r\cdot s$. The empty set and empty word are denoted by $\emptyset$ and $\varepsilon$, respectively.

The language $\mathcal L(r)\subseteq\Sigma^*$ of a standard regular expression $r$ is defined inductively by
\begin{equation}
\begin{aligned}
\mathcal L(\emptyset) &= \emptyset,&
\mathcal L(\varepsilon) &= \{\varepsilon\},&
\mathcal L(a) &= \{a\},\\
\mathcal L(r+s) &= \mathcal L(r)\cup \mathcal L(s),&
\mathcal L(rs) &= \mathcal L(r)\mathcal L(s),&
\mathcal L(r^*) &= \bigcup_{n\ge 0}\mathcal L(r)^n,
\end{aligned}
\end{equation}
where
$\mathcal L(r)\mathcal L(s)
=\{uv \mid u\in\mathcal L(r),\,v\in\mathcal L(s)\}$. We define $\mathsf{Null}(r)\in\{\emptyset,\varepsilon\}$ by $\mathsf{Null}(r)=\varepsilon \iff \varepsilon\in\mathcal{L}(r),$ and $\mathsf{Null}(r)=\emptyset$ otherwise. The \emph{alphabetic size} of $r$, written $|r|_\Sigma$, is the number of alphabet-symbol occurrences in $r$.

% An expression $r$ is \emph{nullable} when it accepts the empty word, namely when $\mathsf{Null}(r)\Longleftrightarrow \varepsilon\in\mathcal L(r)$.


A nondeterministic finite automaton (NFA) is a quintuple $
\mathcal A=\langle Q,\Sigma,\delta,q_0,F\rangle .
\label{eq:nfa-definition}
$ where $Q$ is a finite set of states, $\Sigma$ is a finite
alphabet, $q_0\in Q$ is the initial state, $F\subseteq Q$ is the
set of accepting states, and
$\delta:Q\times\Sigma\longrightarrow 2^Q$ is the transition
function.

A run of $\mathcal A$ on a word $w=a_1a_2\cdots a_n$ is a sequence $
q_0 \xrightarrow{a_1} q_1
\xrightarrow{a_2}\cdots
\xrightarrow{a_n} q_n .
\label{eq:nfa-run}
$ such that $q_t\in\delta(q_{t-1},a_t)$ for every
$1\le t\le n$. The run is accepting when $q_n\in F$. The language
recognized by $\mathcal A$ is
$
\mathcal L(\mathcal A)
=
\{
w\in\Sigma^*
\mid
\mathcal A
\text{ has an accepting run on }w
\}.
\label{eq:nfa-language}
$ An NFA is deterministic if $|\delta(q,a)|\le 1$ for every
$q\in Q$ and $a\in\Sigma$. Two automata are \emph{equivalent}
if they recognize the same language.
\paragraph{Quotient automata.}

Let $\equiv$ be an equivalence relation on $Q$. It is
\emph{right-invariant} with respect to $\mathcal A$ if equivalent
states agree on acceptance and, whenever $p\equiv q$ and
$p'\in\delta(p,a)$, there exists $q'\in\delta(q,a)$ such that
$p'\equiv q'$. The quotient automaton $\mathcal A/{\equiv}$ has equivalence classes
$Q/{\equiv}$ as states and transition relation
$
\delta/{\equiv}([p],a)
=
\{[q]\mid q\in\delta(p,a)\}.
\label{eq:quotient-transition}
$


% ------------------------------------------------------------
\subsection{Derivatives of Regular Expressions}
\label{sec:derivatives-of-regular-expressions}
% ------------------------------------------------------------

For strings $u,v\in\Sigma^*$ with $u\preceq v$, we define the left
quotient of $v$ by $u$ as $
u^{-1}v=z
\:\Longleftrightarrow\:
v=uz.
\label{eq:string-left-quotient}
$
Since $\Sigma^*$ is a free monoid, such a string $z$ is unique. In
particular,
$
\varepsilon^{-1}u=u,
\:
u^{-1}u=\varepsilon.
\label{eq:string-left-quotient-basic}
$
For example,
$
h^{-1}(hello)=ello,
\:
(he)^{-1}(hello)=llo.
\label{eq:string-left-quotient-example}
$

For a language $L\subseteq\Sigma^*$ and a word $w\in\Sigma^*$, the
left quotient of $L$ by $w$ is
$
w^{-1}L
=
\left\{
x\in\Sigma^*
\;\middle|\;
wx\in L
\right\}.
\label{eq:language-left-quotient}
$
Thus, $w^{-1}L$ contains the possible continuations after the prefix
$w$ has been consumed.

The derivative $D_a(r)$ of a regular expression $r$ with respect to a symbol $a\in\Sigma$, introduced by Brzozowski~\cite{brzozowski1964derivatives}, is defined as follows:
\begin{equation}
\begin{aligned}
D_a(\emptyset)
&=\emptyset,
&
D_a(\varepsilon )
&=\emptyset,
&
D_a(b)
&=
\begin{cases}
\varepsilon, & \text{if }a=b,\\
\emptyset, & \text{if }a\neq b,
\end{cases}
\\
D_a(r+s)
&=D_a(r)+D_a(s),
&
D_a(r^*)
&=D_a(r)\cdot r^*,
&
D_a(r\cdot s)
&=
\begin{cases}
D_a(r)\cdot s+D_a(s),
& \text{if }\nu(r),\\
D_a(r)\cdot s,
& \text{otherwise}.
\end{cases}
\end{aligned}
\label{eq:regular-expression-symbol-derivative}
\end{equation}
Here $\nu(r)$ is the nullable predicate, defined by $
\nu(r)
\Longleftrightarrow
\varepsilon\in\mathcal L(r).
\label{eq:regular-expression-nullable}
$ Derivatives with respect to words are defined inductively by
$
D_\varepsilon(r)=r,
\:
D_{aw}(r)=D_w(D_a(r)).
\label{eq:regular-expression-word-derivative}
$
The derivative represents the corresponding language quotient:
$
\mathcal L(D_w(r))
=
w^{-1}\mathcal L(r).
\label{eq:regular-expression-derivative-quotient}
$
Consequently, membership can be decided by the derivative construction
and nullability: $
w\in\mathcal L(r)
\:\Longleftrightarrow\:
\nu(D_w(r)).
\label{eq:regular-expression-derivative-acceptance}
$




% % ------------------------------------------------------------
% \subsection{The Position Automaton}
% \label{sec:position-automaton}
% % ------------------------------------------------------------

% The position automaton, which is constructed by the classical algorithms of Glushkov~\cite{glushkov1961abstract} and of McNaughton and Yamada~\cite{mcnaughton1960regular},
% associates its non-initial states with occurrences of alphabet
% symbols in a regular expression.

% Let $r\in\mathsf{RE}$. Its \emph{marked expression}
% $\overline r$ is obtained by assigning a distinct position to every
% occurrence of an alphabet symbol in $r$. For example,
% $r=(a+b)a$ is marked as
% $\overline r=(a_1+b_2)a_3$.


% Let $\operatorname{Pos}(\overline r)$ denote the set of marked
% positions in $\overline r$. For each
% $p\in\operatorname{Pos}(\overline r)$, let $\sigma_p$ denote the
% marked symbol at position $p$, and let
% $\operatorname{sym}(p)\in\Sigma$ denote its underlying unmarked
% symbol.

% The marked alphabet is
% \begin{equation}
% \overline{\Sigma}_{\overline r}
% =
% \{\sigma_p \mid p\in\operatorname{Pos}(\overline r)\}.
% \label{eq:marked-alphabet}
% \end{equation}


% The classical sets $\operatorname{First}$,
% $\operatorname{Last}$, and $\operatorname{Follow}$ describe the
% relative positions of marked symbol occurrences in words of
% $\mathcal L(\overline r)$. For readability, write
% $\overline\Sigma=\overline{\Sigma}_{\overline r}$. They are defined by
% \begin{equation}
% \begin{gathered}
% \operatorname{First}(\overline r)
% =
% \{p\mid
% \exists u\in\overline\Sigma^{*}.\,
% \sigma_pu\in\mathcal L(\overline r)\},
% \qquad
% \operatorname{Last}(\overline r)
% =
% \{p\mid
% \exists u\in\overline\Sigma^{*}.\,
% u\sigma_p\in\mathcal L(\overline r)\},
% \\[1mm]
% \operatorname{Follow}(\overline r,p)
% =
% \{q\mid
% \exists u,v\in\overline\Sigma^{*}.\,
% u\sigma_p\sigma_qv
% \in\mathcal L(\overline r)\},
% \quad
% p\in\operatorname{Pos}(\overline r).
% \end{gathered}
% \label{eq:first-last-follow-definition}
% \end{equation}


% Thus, in the $\mathcal L(\overline r)$, $\operatorname{First}(\overline r)$ consists of the
% positions that may occur first ,
% $\operatorname{Last}(\overline r)$ consists of the positions that
% may occur last, and $\operatorname{Follow}(\overline r,p)$
% consists of the positions that may occur immediately after $p$.


% Introduce a distinguished initial state $0$, which does not correspond
% to any marked symbol occurrence, and define $\operatorname{Pos}_0(\overline r) = \operatorname{Pos}(\overline r)\cup\{0\}.\label{eq:position-with-initial-state}$ Naturally, we have $\operatorname{Follow}(\overline r,0) = \operatorname{First}(\overline r). \label{eq:initial-follow}$


% The accepting states are
% \begin{equation}
% \operatorname{Last}_0(\overline r)
% =
% \begin{cases}
% \operatorname{Last}(\overline r)\cup\{0\},
% &
% \text{if }\mathsf{Null}(r),
% \\[1mm]
% \operatorname{Last}(\overline r),
% &
% \text{otherwise}.
% \end{cases}
% \label{eq:last-zero-definition}
% \end{equation}

% For $S\subseteq\operatorname{Pos}(\overline r)$ and
% $a\in\Sigma$, define
% \begin{equation}
% \operatorname{Select}(S,a)
% =
% \{p\in S\mid\operatorname{sym}(p)=a\}.
% \label{eq:select-definition}
% \end{equation}


% \begin{definition}[Position automaton]

% The position automaton of $r$ is

% \begin{equation}
% \mathcal A_{\mathrm{pos}}(r)
% =
% \left\langle
% \operatorname{Pos}_0(\overline r),
% \Sigma,
% \delta_{\mathrm{pos}},
% 0,
% \operatorname{Last}_0(\overline r)
% \right\rangle .
% \label{eq:position-automaton}
% \end{equation}

% \end{definition}


% Hence, $\operatorname{First}(\overline r)$ determines the states
% reachable from the initial state,
% $\operatorname{Follow}(\overline r,p)$ determines the successors of a
% position $p$, and $\operatorname{Last}_0(\overline r)$ determines the
% accepting states.

% Every transition enters a marked symbol occurrence and consumes its
% corresponding input symbol. Consequently, the classical position
% automaton is $\varepsilon$-free. Its standard correctness property is
% \begin{equation}
% \mathcal L(\mathcal A_{\mathrm{pos}}(r))
% =
% \mathcal L(r).
% \label{eq:position-correctness}
% \end{equation}


